% Part: first-order-logic
% Chapter: introduction
% Section: semantic-notions

\documentclass[../../../include/open-logic-section]{subfiles}

\begin{document}

\olfileid[pt]{fol}{int}{sem}

\olsection{Noções Semânticas}

Mencionamos acima que, quando consideramos se $\Sat{M}{!A}[s]$ vale,
(por conveniência) deixamos $s$ atribuir valores a todas as
!!{variable}s, mas apenas os valores que atribui às !!{variable}s
em~$!A$ são usados.  De fato, são apenas os valores das
!!{variable}s \emph{livres} em~$!A$ que importam. É claro que, por
sermos cuidadosos, vamos provar esse fato. Como as !!{sentence}s não têm
variáveis livres, $s$~não importa de modo algum quando se trata de
saber se são satisfeitas em !!a{structure}.  Assim, quando $!A$ é
!!a{sentence}, podemos definir $\Sat{M}{!A}$ como significando
``$\Sat{M}{!A}[s]$ para todo~$s$,'' o que, aliás, é verdadeiro se, e
somente se, $\Sat{M}{!A}[s]$ para pelo menos um~$s$. Precisamos
introduzir atribuições de variáveis para obter uma definição funcional
de satisfação para !!{formula}s, mas para !!{sentence}s, a satisfação é
independente das atribuições de variáveis.

Uma vez que temos uma definição de ``$\Sat{M}{!A}$,'' sabemos o que
``caso'' e ``verdadeira em'' significam no que diz respeito às
!!{sentence}s da lógica de primeira ordem. Com base na definição de
$\Sat{M}{!A}$ para !!{sentence}s, podemos então definir as noções
semânticas básicas de validade, consequência e satisfatibilidade.  Uma
sentença é válida, $\Entails !A$, se toda !!{structure} a satisfaz. Ela
é consequência de um conjunto de !!{sentence}s, $\Gamma \Entails !A$,
se toda !!{structure} que satisfaz todas as !!{sentence}s em~$\Gamma$
também satisfaz~$!A$. E um conjunto de !!{sentence}s é satisfatível se
alguma !!{structure} satisfaz todas as !!{sentence}s nele ao mesmo
tempo.

Como as !!{formula}s são definidas indutivamente, e a satisfação é, por
sua vez, definida por indução na estrutura das !!{formula}s, podemos
usar indução para provar propriedades da nossa semântica e relacionar as
noções semânticas definidas.  Vamos reunir e provar algumas dessas
propriedades, em parte porque são individualmente interessantes, mas
principalmente porque muitas delas serão úteis quando formos investigar
a relação entre semântica e sistemas de !!{derivation}. Para isso,
também teremos que definir (precisamente, isto é, por indução) algumas
noções sintáticas e operações que ainda não mencionamos.

\end{document}
